\subsection{Syntaktická analýza}
\label{sec:parser}

Parser (\texttt{src/parser.rs}) je implementován jako jednoduchý
rekurzivně-sestupný analyzátor nad vektorem tokenů. Vzhledem
k~tomu, že zásobníkový jazyk má lineární zápis (program je sekvence slov
a~literálů), je výsledný abstraktní syntaktický strom (AST) záměrně mělký.

Kořenem AST je uzel \texttt{Program}, který obsahuje seznam příkazů. Jedinou
hierarchickou konstrukcí jsou definice slov uzavřené mezi \texttt{:} a~\texttt{;},
které jsou reprezentovány uzlem \texttt{Definition} s~vlastním tělem.

Základní uzly AST jsou definovány v~\texttt{src/types.rs}.

\begin{lstlisting}[language=Rust, caption={Základní uzly AST}]
pub enum AstNode {
    Number(i32, Position),
    Word(String, Position),
    StringLiteral(String, Position),
    Definition { name: String, body: Vec<AstNode>, position: Position },
    VariableDeclaration { name: String, position: Position },
    Program(Vec<AstNode>),
}
\end{lstlisting}

Parser v~celém programu ignoruje tokeny typu \texttt{Comment}. Tím se snižuje
komplexita syntaktické analýzy bez ztráty informace o~pozicích chyb, protože
pozice je nesena i~ostatními tokeny.

Definice slova je parsována následovně: po \texttt{:} musí následovat jméno
slova, poté parser postupně načítá uzly těla až do \texttt{;}. Vnořené definice
jsou zakázány a~jsou vyhodnoceny jako syntaktická chyba.

Parser zároveň implementuje speciální pravidlo pro konstrukci \texttt{VARIABLE}.
Pokud se v~toku tokenů objeví slovo \texttt{VARIABLE}, parser očekává následující
identifikátor, který uloží do uzlu \texttt{VariableDeclaration}. Vlastní práce
s~proměnnými je řešena až v~nižších fázích (sémantika a~IR).
